\IfFileExists{stacks-project.cls}{%
\documentclass{stacks-project}
}{%
\documentclass{amsart}
}

% For dealing with references we use the comment environment
\usepackage{verbatim}
\newenvironment{reference}{\comment}{\endcomment}
%\newenvironment{reference}{}{}
\newenvironment{slogan}{\comment}{\endcomment}
\newenvironment{history}{\comment}{\endcomment}

% For commutative diagrams we use Xy-pic
\usepackage[all]{xy}

% We use 2cell for 2-commutative diagrams.
\xyoption{2cell}
\UseAllTwocells

% We use multicol for the list of chapters between chapters
\usepackage{multicol}

% This is generally recommended for better output
\usepackage{lmodern}
\usepackage[T1]{fontenc}

% For cross-file-references
\usepackage{xr-hyper}

% Package for hypertext links:
\usepackage{hyperref}

% For any local file, say "hello.tex" you want to link to please
% use \externaldocument[hello-]{hello}
\externaldocument[introduction-]{introduction}
\externaldocument[conventions-]{conventions}
\externaldocument[sets-]{sets}
\externaldocument[categories-]{categories}
\externaldocument[topology-]{topology}
\externaldocument[sheaves-]{sheaves}
\externaldocument[sites-]{sites}
\externaldocument[stacks-]{stacks}
\externaldocument[fields-]{fields}
\externaldocument[algebra-]{algebra}
\externaldocument[brauer-]{brauer}
\externaldocument[homology-]{homology}
\externaldocument[derived-]{derived}
\externaldocument[simplicial-]{simplicial}
\externaldocument[more-algebra-]{more-algebra}
\externaldocument[smoothing-]{smoothing}
\externaldocument[modules-]{modules}
\externaldocument[sites-modules-]{sites-modules}
\externaldocument[injectives-]{injectives}
\externaldocument[cohomology-]{cohomology}
\externaldocument[sites-cohomology-]{sites-cohomology}
\externaldocument[dga-]{dga}
\externaldocument[dpa-]{dpa}
\externaldocument[sdga-]{sdga}
\externaldocument[hypercovering-]{hypercovering}
\externaldocument[schemes-]{schemes}
\externaldocument[constructions-]{constructions}
\externaldocument[properties-]{properties}
\externaldocument[morphisms-]{morphisms}
\externaldocument[coherent-]{coherent}
\externaldocument[divisors-]{divisors}
\externaldocument[limits-]{limits}
\externaldocument[varieties-]{varieties}
\externaldocument[topologies-]{topologies}
\externaldocument[descent-]{descent}
\externaldocument[perfect-]{perfect}
\externaldocument[more-morphisms-]{more-morphisms}
\externaldocument[flat-]{flat}
\externaldocument[groupoids-]{groupoids}
\externaldocument[more-groupoids-]{more-groupoids}
\externaldocument[etale-]{etale}
\externaldocument[chow-]{chow}
\externaldocument[intersection-]{intersection}
\externaldocument[pic-]{pic}
\externaldocument[weil-]{weil}
\externaldocument[adequate-]{adequate}
\externaldocument[dualizing-]{dualizing}
\externaldocument[duality-]{duality}
\externaldocument[discriminant-]{discriminant}
\externaldocument[derham-]{derham}
\externaldocument[local-cohomology-]{local-cohomology}
\externaldocument[algebraization-]{algebraization}
\externaldocument[curves-]{curves}
\externaldocument[resolve-]{resolve}
\externaldocument[models-]{models}
\externaldocument[functors-]{functors}
\externaldocument[equiv-]{equiv}
\externaldocument[pione-]{pione}
\externaldocument[etale-cohomology-]{etale-cohomology}
\externaldocument[proetale-]{proetale}
\externaldocument[relative-cycles-]{relative-cycles}
\externaldocument[more-etale-]{more-etale}
\externaldocument[trace-]{trace}
\externaldocument[crystalline-]{crystalline}
\externaldocument[spaces-]{spaces}
\externaldocument[spaces-properties-]{spaces-properties}
\externaldocument[spaces-morphisms-]{spaces-morphisms}
\externaldocument[decent-spaces-]{decent-spaces}
\externaldocument[spaces-cohomology-]{spaces-cohomology}
\externaldocument[spaces-limits-]{spaces-limits}
\externaldocument[spaces-divisors-]{spaces-divisors}
\externaldocument[spaces-over-fields-]{spaces-over-fields}
\externaldocument[spaces-topologies-]{spaces-topologies}
\externaldocument[spaces-descent-]{spaces-descent}
\externaldocument[spaces-perfect-]{spaces-perfect}
\externaldocument[spaces-more-morphisms-]{spaces-more-morphisms}
\externaldocument[spaces-flat-]{spaces-flat}
\externaldocument[spaces-groupoids-]{spaces-groupoids}
\externaldocument[spaces-more-groupoids-]{spaces-more-groupoids}
\externaldocument[bootstrap-]{bootstrap}
\externaldocument[spaces-pushouts-]{spaces-pushouts}
\externaldocument[spaces-chow-]{spaces-chow}
\externaldocument[groupoids-quotients-]{groupoids-quotients}
\externaldocument[spaces-more-cohomology-]{spaces-more-cohomology}
\externaldocument[spaces-simplicial-]{spaces-simplicial}
\externaldocument[spaces-duality-]{spaces-duality}
\externaldocument[formal-spaces-]{formal-spaces}
\externaldocument[restricted-]{restricted}
\externaldocument[spaces-resolve-]{spaces-resolve}
\externaldocument[formal-defos-]{formal-defos}
\externaldocument[defos-]{defos}
\externaldocument[cotangent-]{cotangent}
\externaldocument[examples-defos-]{examples-defos}
\externaldocument[algebraic-]{algebraic}
\externaldocument[examples-stacks-]{examples-stacks}
\externaldocument[stacks-sheaves-]{stacks-sheaves}
\externaldocument[criteria-]{criteria}
\externaldocument[artin-]{artin}
\externaldocument[quot-]{quot}
\externaldocument[stacks-properties-]{stacks-properties}
\externaldocument[stacks-morphisms-]{stacks-morphisms}
\externaldocument[stacks-limits-]{stacks-limits}
\externaldocument[stacks-cohomology-]{stacks-cohomology}
\externaldocument[stacks-perfect-]{stacks-perfect}
\externaldocument[stacks-introduction-]{stacks-introduction}
\externaldocument[stacks-more-morphisms-]{stacks-more-morphisms}
\externaldocument[stacks-geometry-]{stacks-geometry}
\externaldocument[moduli-]{moduli}
\externaldocument[moduli-curves-]{moduli-curves}
\externaldocument[examples-]{examples}
\externaldocument[exercises-]{exercises}
\externaldocument[guide-]{guide}
\externaldocument[desirables-]{desirables}
\externaldocument[coding-]{coding}
\externaldocument[obsolete-]{obsolete}
\externaldocument[fdl-]{fdl}
\externaldocument[index-]{index}

% Theorem environments.
%
\theoremstyle{plain}
\newtheorem{theorem}[subsection]{Theorem}
\newtheorem{proposition}[subsection]{Proposition}
\newtheorem{lemma}[subsection]{Lemma}

\theoremstyle{definition}
\newtheorem{definition}[subsection]{Definition}
\newtheorem{example}[subsection]{Example}
\newtheorem{exercise}[subsection]{Exercise}
\newtheorem{situation}[subsection]{Situation}

\theoremstyle{remark}
\newtheorem{remark}[subsection]{Remark}
\newtheorem{remarks}[subsection]{Remarks}

\numberwithin{equation}{subsection}

% Macros
%
\def\lim{\mathop{\mathrm{lim}}\nolimits}
\def\colim{\mathop{\mathrm{colim}}\nolimits}
\def\Spec{\mathop{\mathrm{Spec}}}
\def\Hom{\mathop{\mathrm{Hom}}\nolimits}
\def\Ext{\mathop{\mathrm{Ext}}\nolimits}
\def\SheafHom{\mathop{\mathcal{H}\!\mathit{om}}\nolimits}
\def\SheafExt{\mathop{\mathcal{E}\!\mathit{xt}}\nolimits}
\def\Sch{\mathit{Sch}}
\def\Mor{\mathop{\mathrm{Mor}}\nolimits}
\def\Ob{\mathop{\mathrm{Ob}}\nolimits}
\def\Sh{\mathop{\mathit{Sh}}\nolimits}
\def\NL{\mathop{N\!L}\nolimits}
\def\CH{\mathop{\mathrm{CH}}\nolimits}
\def\proetale{{pro\text{-}\acute{e}tale}}
\def\etale{{\acute{e}tale}}
\def\QCoh{\mathit{QCoh}}
\def\Ker{\mathop{\mathrm{Ker}}}
\def\Im{\mathop{\mathrm{Im}}}
\def\Coker{\mathop{\mathrm{Coker}}}
\def\Coim{\mathop{\mathrm{Coim}}}

% Boxtimes
%
\DeclareMathSymbol{\boxtimes}{\mathbin}{AMSa}{"02}

%
% Macros for moduli stacks/spaces
%
\def\QCohstack{\mathcal{QC}\!\mathit{oh}}
\def\Cohstack{\mathcal{C}\!\mathit{oh}}
\def\Spacesstack{\mathcal{S}\!\mathit{paces}}
\def\Quotfunctor{\mathrm{Quot}}
\def\Hilbfunctor{\mathrm{Hilb}}
\def\Curvesstack{\mathcal{C}\!\mathit{urves}}
\def\Polarizedstack{\mathcal{P}\!\mathit{olarized}}
\def\Complexesstack{\mathcal{C}\!\mathit{omplexes}}
% \Pic is the operator that assigns to X its picard group, usage \Pic(X)
% \Picardstack_{X/B} denotes the Picard stack of X over B
% \Picardfunctor_{X/B} denotes the Picard functor of X over B
\def\Pic{\mathop{\mathrm{Pic}}\nolimits}
\def\Picardstack{\mathcal{P}\!\mathit{ic}}
\def\Picardfunctor{\mathrm{Pic}}
\def\Deformationcategory{\mathcal{D}\!\mathit{ef}}
\title{The triangle attached to an admissible short exact sequence is independent of the choices made}
\author{The Stacks Project Authors}
\date{stacks-project:dga.tex:local-snapshot-20261012; GNU FDL 1.2 or later}
\begin{document}
\maketitle
\noindent\textit{Editorial note.} The counted statement is the author's own lemma environment labelled lemma-analogue-triangle-independent-splittings (source lines 4104-4115); the proof environment begins two lines below it and is its own proof, closing at source line 4142.\par\medskip
\noindent\textit{Editorial note.} The statement and the proof invoke the standing situation and the earlier construction lemma by label; both are retained as context spans, and the generator expands the context transitively, so the references resolve inside the assembled package.\par\medskip
\noindent\textit{Editorial note.} The statement displays one triangle diagram; it is part of the retained source range and is reproduced as written, so no figure file is needed.\par\medskip
\noindent\textit{Editorial note.} The retained preamble is the chapter's own minimal preamble (sources/stacks-dga/preamble.tex, taken with its real line count), a verbatim excerpt of the native Xy-pic setup, the native theorem declarations with the mathematical macros, and the project's own no-op reference environments; no macro was invented or redefined.\par\medskip
\noindent\textit{Editorial scope.} One original result of the Stacks Project's Differential Graded Algebra chapter, counted once: the triangle attached to an admissible short exact sequence is independent of the choices made in its construction, up to canonical isomorphism.\par\medskip
\noindent\textit{Allowed standard prerequisites (editorial).}\begin{itemize}
\item The standing situation (axioms A, B, C). A dg category with the distinguished-triangle axioms, including the construction of the connecting morphism of an admissible short exact sequence (the author's Situation situation-ABC, retained as context).. It supplies the hypotheses and the structure maps of the two constructions being compared (source lines 4106-4114, 4117-4130).
\item The triangle attached to an admissible sequence. The construction of the triangle with its connecting morphism delta from a choice of splitting (the author's lemma lemma-get-triangle, retained as context).. It is the construction whose independence is asserted, and formula (5) of it is used in the verification (source lines 4112-4114, 4140-4141).
\item Homotopy of the connecting morphisms. Two connecting morphisms built from different splittings differ by an explicit homotopy, so the resulting triangles are canonically isomorphic in the homotopy category.. It is the content of the proof (source lines 4117-4142).
\end{itemize}\medskip
\begin{situation}
\label{situation-ABC}
Here $R$ is a ring and $\mathcal{A}$ is a differential graded category
over $R$ having axioms (A), (B), and
\begin{enumerate}
\item[(C)] given an arrow $f : x \to y$ of degree $0$ with
$\text{d}(f) = 0$ there exists an admissible short exact sequence
$y \to c(f) \to x[1]$ in $\text{Comp}(\mathcal{A})$ such that the map
$x[1] \to y[1]$ of Lemma \ref{lemma-get-triangle} is equal to $f[1]$.
\end{enumerate}
\end{situation}

\begin{lemma}
\label{lemma-get-triangle}
Let $\mathcal{A}$ be a differential graded category satisfying
axioms (A) and (B). Given an admissible short exact sequence
$x \to y \to z$ we obtain (see proof) a triangle
$$
x \to y \to z \to x[1]
$$
in $\text{Comp}(\mathcal{A})$ with the property that any two compositions
in $z[-1] \to x \to y \to z \to x[1]$ are zero in $K(\mathcal{A})$.
\end{lemma}

\begin{proof}
Choose a diagram
$$
\xymatrix{
x \ar[rr]_1 \ar[rd]_a & & x \\
& y \ar[ru]_\pi \ar[rd]^b & \\
z \ar[rr]^1 \ar[ru]^s & & z
}
$$
giving the isomorphism of graded objects $y \cong x \oplus z$ as in the
definition of an admissible short exact sequence. Here are some equations
that hold in this situation
\begin{enumerate}
\item $1 = \pi a$ and hence $\text{d}(\pi) a = 0$,
\item $1 = b s$ and hence $b \text{d}(s) = 0$,
\item $1 = a \pi + s b$ and hence $a \text{d}(\pi) + \text{d}(s) b = 0$,
\item $\pi s = 0$ and hence $\text{d}(\pi)s + \pi \text{d}(s) = 0$,
\item $\text{d}(s) = a \pi \text{d}(s)$ because
$\text{d}(s) = (a \pi + s b)\text{d}(s)$ and $b\text{d}(s) = 0$,
\item $\text{d}(\pi) = \text{d}(\pi) s b$ because
$\text{d}(\pi) = \text{d}(\pi)(a \pi + s b)$ and $\text{d}(\pi)a = 0$,
\item $\text{d}(\pi \text{d}(s)) = 0$ because if we postcompose it
with the monomorphism $a$ we get
$\text{d}(a\pi \text{d}(s)) = \text{d}(\text{d}(s)) = 0$, and
\item $\text{d}(\text{d}(\pi)s) = 0$ as by (4) it is the negative
of $\text{d}(\pi\text{d}(s))$ which is $0$ by (7).
\end{enumerate}
We've used repeatedly that $\text{d}(a) = 0$, $\text{d}(b) = 0$,
and that $\text{d}(1) = 0$. By (7) we see that
$$
\delta = \pi \text{d}(s) = - \text{d}(\pi) s : z \to x[1]
$$
is a morphism in $\text{Comp}(\mathcal{A})$. By (5) we see that
the composition $a \delta = a \pi \text{d}(s) = \text{d}(s)$
is homotopic to zero. By (6) we see that the composition
$\delta b = - \text{d}(\pi)sb = \text{d}(-\pi)$ is homotopic to zero.
\end{proof}

\begin{lemma}
\label{lemma-analogue-triangle-independent-splittings}
In Situation \ref{situation-ABC}
let $0 \to x \to y \to z \to 0$ be an admissible short
exact sequence in $\text{Comp}(\mathcal{A})$. The triangle
$$
\xymatrix{x\ar[r] & y\ar[r] & z\ar[r]^{\delta} & x[1]}
$$
with $\delta : z \to x[1]$ as defined in Lemma \ref{lemma-get-triangle}
is up to canonical isomorphism in $K(\mathcal{A})$, independent of the
choices made in Lemma \ref{lemma-get-triangle}.
\end{lemma}

\begin{proof}
Suppose $\delta$ is defined by the splitting
$$
\xymatrix{
x \ar@<0.5ex>[r]^{a} &
y \ar@<0.5ex>[r]^b\ar@<0.5ex>[l]^{\pi} &
z \ar@<0.5ex>[l]^s
}
$$
and $\delta'$ is defined by the splitting with $\pi',s'$
in place of $\pi,s$. Then
$$
s'-s = (a\pi + sb)(s'-s) = a\pi s'
$$
since $bs' = bs = 1_z$ and $\pi s = 0$. Similarly,
$$
\pi' - \pi = (\pi' - \pi)(a\pi + sb) = \pi'sb
$$
Since $\delta = \pi d(s)$ and $\delta' = \pi'd(s')$
as constructed in Lemma \ref{lemma-get-triangle}, we may compute
$$
\delta' = \pi'd(s') = (\pi + \pi'sb)d(s + a\pi s') = \delta + d(\pi s')
$$
using $\pi a = 1_x$, $ba = 0$, and $\pi'sbd(s') = \pi'sba\pi d(s') = 0$
by formula (5) in Lemma \ref{lemma-get-triangle}.
\end{proof}

\end{document}
